\section{What Happens in the Operating Room}

\subsection{The room, the cabinet, and the cable that is not connected}

Three things cross the network boundary of the operating suite before a
procedure: a cryptographically signed model build, the protocol version, and the
pinned registry entry that binds the two. Two things cross afterwards: the
required reports to the FDA and to the reviewing Institutional Review Board.
During the procedure itself nothing crosses, because during the procedure no
route exists.

That is the whole answer to the cybersecurity concern, and it is architectural
rather than reassuring. If there is no path out of the building while the
operation is running, the attack surface during the operation is the building.
The relevant FDA guidance frames the expectation for premarket cybersecurity
documentation \cite{FDACyber2026}; the design that satisfies it here is
on-premises inference with the external interface closed for the duration
\cite{onpremwhippl}.

Figure 17 draws the stack. The topology is adapted from the author's autonomous
single-patient journey simulation \cite{paipatientjr}, which was non-small cell
lung cancer, and it is re-scoped here to PDAC. Four differences matter: eight
arms rather than four, three anastomoses rather than one bronchial closure, a
vascular exclusion envelope around the superior mesenteric and portal veins that
the thoracic configuration has no equivalent of, and a drain that stays in for
days rather than hours.

\begin{pafloat}
\begin{pafig}[diagrams (python)-type: deployment, full page]
\dgnodew{dgtiled}{-6.4}{0}{\glyphrobot}{8 robotic arms, $\leq 3$ N each, $\leq 18$ N summed}{r1}
\dgnode{dgtile}{-3.6}{0}{\glyphsignal}{Sensor stack, 640 channels, 10 kHz}{r2}
\dgnodew{dgtiled}{-0.8}{0}{\glyphstop}{Emergency stop, $\leq 3$ ms, $\leq 500$ ms}{r3}
\dgnode{dgtile}{-6.4}{-3.0}{\glyphmon}{Surgeon console, sole approval authority}{r4}
\dgnode{dgtile}{-3.6}{-3.0}{\glyphmon}{Second operator station, independent stop}{r5}
\dgnode{dgtile}{-0.8}{-3.0}{\glyphuser}{Independent safety monitor, every case}{r6}
\dgnode{dgtile}{4.6}{0}{\glyphai}{On-premises LLM, advisory only, no actuator path}{c1}
\dgnodew{dgtiled}{7.4}{0}{\glyphshield}{Deterministic gate, $\leq 2$ mm, $\leq 0.5$ N}{c2}
\dgnode{dgtile}{10.2}{0}{\glyphdb}{Version registry, frozen at your consent}{c3}
\dgnode{dgtile}{4.6}{-3.0}{\glyphlock}{Hash-chained audit trail, 21 CFR part 11}{c4}
\dgnode{dgtile}{7.4}{-3.0}{\glyphdb}{Data vault, access list published}{c5}
\dgnode{dgtile}{10.2}{-3.0}{\glyphdoc}{Your day-30 and day-90 results}{c6}
\dgnodeg{dgtilek}{1.0}{-6.6}{\glyphnet}{Network boundary, closed for the whole procedure}{b1}
\dgnodeg{dgtileg}{4.4}{-6.6}{\glyphcloud}{Hospital network}{b2}
\dgnodeg{dgtileg}{7.8}{-6.6}{\glyphlink}{Vendor support and remote assistance}{b3}
\begin{scope}[on background layer]
\node[dgcluster,fit=(r1)(r1l)(r2)(r2l)(r3)(r3l)] (K1) {};
\node[dgcluster2,fit=(r4)(r4l)(r5)(r5l)(r6)(r6l)] (K2) {};
\node[dgcluster,fit=(c1)(c1l)(c2)(c2l)(c3)(c3l)] (K3) {};
\node[dgcluster2,fit=(c4)(c4l)(c5)(c5l)(c6)(c6l)] (K4) {};
\node[dgcluster2,fit=(b1)(b1l)(b2)(b2l)(b3)(b3l)] (K5) {};
\end{scope}
\node[dgctitle,anchor=south west]  at (K1.north west) {sterile field, the table you are on};
\node[dgctitle2,anchor=south west] at (K2.north west) {operating room, non-sterile side};
\node[dgctitle,anchor=south west]  at (K3.north west) {on-premises cabinet, same building};
\node[dgctitle2,anchor=south west] at (K4.north west) {evidence, written once and never edited};
\node[dgctitle2,anchor=south west] at (K5.north west) {outside the building};
\draw[dgedge]  (r2) -- node[font=\tiny\sffamily,above]{field state} (r1);
\draw[dgedgek] (r3) to[out=180,in=-20,looseness=0.85] node[font=\tiny\sffamily,below,pos=0.4]{halt, unconditional} (r1);
\draw[dgedgeb] (r4) -- node[font=\tiny\sffamily,left,pos=0.5]{signed motion} (r1);
\draw[dgedgek] (r5) to[out=60,in=-120,looseness=0.85] node[font=\tiny\sffamily,left,pos=0.5]{may stop} (r3);
\draw[dgedgek] (r6) -- node[font=\tiny\sffamily,right,pos=0.5]{may stop} (r3);
\draw[dgedge]  (r2) to[out=20,in=160,looseness=0.85] (c1);
\draw[dgedged] (c1) -- node[font=\tiny\sffamily,above]{candidate plan} (c2);
\draw[dgedgeb] (c2) to[out=-160,in=20,looseness=0.85] node[font=\tiny\sffamily,above,pos=0.45]{only passed plans} (r4);
\draw[dgedge]  (c3) to[out=-160,in=20,looseness=0.85] node[font=\tiny\sffamily,above,pos=0.4]{pinned build} (c1);
\draw[dgedge]  (c2) -- (c4);
\draw[dgedge]  (c4) -- (c5);
\draw[dgedgeb] (c5) -- (c6);
\draw[dgedged] (b1) -- node[font=\tiny\sffamily,above]{no route} (b2);
\draw[dgedged] (b2) -- (b3);
\draw[dgedged] (b1) to[out=90,in=-135,looseness=0.85] node[font=\tiny\sffamily,left,pos=0.55]{before: signed model build} (c1.south west);
\draw[dgedged] (b1) to[out=45,in=-100,looseness=0.85] node[font=\tiny\sffamily,right,pos=0.6]{before: protocol version} (c2.south);
\draw[dgedged] (c4) to[out=-135,in=60,looseness=0.85] node[font=\tiny\sffamily,left,pos=0.4]{after: reports to FDA and IRB} (b1.north east);
\draw[dgedgek] (c1) to[out=-110,in=110,looseness=0.9] (b1.north);
\pxmark{2.5}{-4.5}
\node[font=\tiny\sffamily\bfseries,fill=protowhite,inner sep=1.6pt,anchor=west] at (2.75,-4.5) {DURING the procedure: no route exists};
\node[font=\scriptsize\sffamily\itshape,text=protogray,anchor=north west,align=left]
  at (-8.2,-8.6) {Three things cross the boundary before a procedure and two cross afterwards. Nothing crosses during it. The struck line is the\\claim: on-premises inference with the external interface closed means the attack surface during the operation is the building.};
\end{pafig}
\vspace{-0.7cm}
\figcaption{Figure 17. The operating room and the on-premises stack drawn as a deployment architecture,\\
with the sterile field, the non-sterile side, the cabinet, the evidence store, and the\\
network boundary that carries no intra-procedural route out of the building at all.}
\end{pafloat}

\subsection{The eight arms, and what limits them}

The limits below are enforced beneath the level of the software that makes the
decision. A model that proposed an out-of-envelope motion would have that
proposal rejected by a deterministic gate; a controller that commanded an
out-of-envelope force would have the arm halted by the force sensing itself. The
distinction matters, because a limit implemented in the same layer as the
decision is not a limit, it is a preference.

\vspace{0.30em}

\begin{xltabular}{\textwidth}{L{4.4cm} C{2.3cm} C{2.4cm} Y}
\toprule
\textbf{Limit} & \textbf{Protocol cap} & \textbf{Worst observed in simulation} & \textbf{What it protects} \\
\midrule
\endfirsthead
\multicolumn{4}{@{}l}{\footnotesize\itshape Table continued from the previous page.}\\
\toprule
\textbf{Limit} & \textbf{Protocol cap} & \textbf{Worst observed in simulation} & \textbf{What it protects} \\
\midrule
\endhead
\midrule
\multicolumn{4}{r@{}}{\footnotesize\itshape Table continued on the next page.}\\
\endfoot
\bottomrule
\endlastfoot
Tip force, any single arm & 3.0 N & 2.1 N & Tissue crush injury at a single instrument \\
Tip force, all arms summed & 18.0 N & 12.4 N & Cumulative traction on the mesentery and the vessels \\
Positional tolerance & 2.0 mm & 1.3 mm & Margin accuracy and corridor integrity \\
Force reproducibility & 0.5 N & 0.31 N & Consistency of the anastomotic suture line \\
Heartbeat period & 0.1 ms & 0.1 ms, held & Detection of any loss of control before motion continues \\
Cross-arm emergency stop & 3.0 ms & 1.7 ms & Time from the surgeon pressing stop to all arms held \\
System-wide emergency stop & 500 ms & 310 ms & Time to a fully quiescent, safe-pose system \\
No-fly corridor & 4 mm either side & no entries & Superior mesenteric vein and portal vein \\
\end{xltabular}

\vspace{0.30em}

\noindent The observed column is class S: it comes from the author's simulation
set \cite{pdac060s2030} and not from human procedures. A simulated worst case is
not a human worst case, and \S 10 repeats that on the same line as the numbers.
The standard the architecture is built against is the international standard for
robotically assisted surgical equipment \cite{iec8060127}.

\subsection{Who is controlling the operation}

This is the most frequently raised of the twenty-one concerns and the one a
paragraph answers worst. The chain has four links and one absence.

The sensor stack streams the current field state: force vectors, vessel
proximity, instrument pose, across 640 channels at 10 kHz. The on-premises model
reads that stream and emits a candidate motion plan with a confidence value and
a stated rationale. The deterministic gate tests the candidate against every
limit in the table above; if any limit would be violated, the plan is rejected,
the rejection is logged with the violated limit named, and the plan is never
presented to the surgeon as executable. If every limit is satisfied, the plan is
presented for approval, with its confidence and rationale attached. The
operating surgeon accepts, modifies, or declines it. Declining requires no
justification. Only an approved, surgeon-signed command reaches the arms.

The absence is that there is no message from the model to the arms. Not a
throttled one, not a supervised one, not a fallback one. Figure 19 draws five
lifelines and the arrow that would be needed is not there. That is what a
systematic review of FDA-cleared surgical robots describes as task-level
autonomy under supervision, the highest level any cleared system reaches
\cite{LeeAuto2024X}, and it is what the FDA's own description of
computer-assisted surgical systems assumes \cite{FDARobot2022}.

Deference is handled separately, because the architecture does not address it.
If a surgeon always accepts, the human in the loop is a formality. The protocol
therefore records every override with its reason and publishes the override rate
as an exploratory endpoint, so automation bias is a measured quantity rather
than an assumed absence \cite{FDATrans2024}.

\subsection{How fast stopping actually is}

Three budgets are specified, and Figure 18 puts all three on one logarithmic
millisecond axis because they span four orders of magnitude and a linear axis
would render the first two invisible.

The heartbeat bus broadcasts at 10 kHz, a 0.1 ms period. Loss of the heartbeat
halts every arm without any human involvement. When a human does press stop, the
cross-arm budget is 3 ms from assertion to all arms held, and the system-wide
budget is 500 ms to a fully quiescent state. The gap between the two is
deliberate margin, not unaccounted latency.

For scale: a human blink takes roughly 100 to 150 ms. Every arm is held before a
blink completes, and the system-wide guarantee is under four blinks. Three
parties may press stop: the operating surgeon at the console, the second
operator at an independent station, and the independent safety monitor who is
present for every procedure. The fail-safe design standard for autonomous and
semi-autonomous systems frames what such a guarantee has to demonstrate
\cite{ieee7009}.

\begin{pafloat}
\begin{pafig}[plantuml-type: timing diagram, @startuml]
\foreach \r/\lab in {0/{Surgeon intent},1/{Heartbeat bus, 10 kHz},2/{Arm 1 to 8 motion},3/{Force at the tip},4/{Protocol budget line}}{
  \node[font=\tiny\sffamily,anchor=east,text width=30mm,align=right] at (-0.3,{-\r*1.55}) {\lab};
  \draw[protogray,line width=0.3pt] (0,{-\r*1.55}) -- (14.0,{-\r*1.55});}
\foreach \x/\lab in {0/{0.1},3.5/{1},7.0/{10},10.5/{100},14.0/{1000}}{
  \draw[protogray,line width=0.4pt] (\x,-6.5) -- (\x,-6.3);
  \node[font=\tiny\sffamily,anchor=north] at (\x,-6.5) {\lab};}
\draw[protogray,line width=0.4pt] (0,-6.3) -- (14.0,-6.3);
\node[font=\tiny\sffamily,anchor=north] at (7.0,-7.0) {elapsed time, milliseconds, logarithmic};
\fill[pagraym,draw=protoblack,line width=0.3pt] (0,-0.2) rectangle (3.5,0.2);
\node[font=\tiny\sffamily] at (1.75,0) {Operating};
\fill[protoblue,draw=protoblack,line width=0.3pt] (3.5,-0.2) rectangle (12.6,0.2);
\node[font=\tiny\sffamily,text=protowhite] at (8.0,0) {Stop pressed};
\fill[pagraym,draw=protoblack,line width=0.3pt] (12.6,-0.2) rectangle (14.0,0.2);
\node[font=\tiny\sffamily] at (13.3,0) {Operating};
\fill[pagraym,draw=protoblack,line width=0.3pt] (0,-1.75) rectangle (3.9,-1.35);
\node[font=\tiny\sffamily] at (1.95,-1.55) {Nominal};
\fill[protoblue,draw=protoblack,line width=0.3pt] (3.9,-1.75) rectangle (6.8,-1.35);
\node[font=\tiny\sffamily,text=protowhite] at (5.35,-1.55) {Halt broadcast};
\fill[pablue2,draw=protoblack,line width=0.3pt] (6.8,-1.75) rectangle (14.0,-1.35);
\node[font=\tiny\sffamily] at (10.4,-1.55) {Quiescent};
\fill[pagraym,draw=protoblack,line width=0.3pt] (0,-3.30) rectangle (5.9,-2.90);
\node[font=\tiny\sffamily] at (2.95,-3.10) {Commanded};
\fill[protoblue,draw=protoblack,line width=0.3pt] (5.9,-3.30) rectangle (6.8,-2.90);
\node[font=\tiny\sffamily,text=protowhite,anchor=west] at (6.95,-3.55) {Decelerating};
\fill[pablue2,draw=protoblack,line width=0.3pt] (6.8,-3.30) rectangle (14.0,-2.90);
\node[font=\tiny\sffamily] at (10.4,-3.10) {Held};
\fill[pagraym,draw=protoblack,line width=0.3pt] (0,-4.85) rectangle (6.6,-4.45);
\node[font=\tiny\sffamily] at (3.3,-4.65) {Within cap};
\fill[protoblue,draw=protoblack,line width=0.3pt] (6.6,-4.85) rectangle (7.6,-4.45);
\fill[pablue2,draw=protoblack,line width=0.3pt] (7.6,-4.85) rectangle (14.0,-4.45);
\node[font=\tiny\sffamily] at (10.8,-4.65) {Zero};
\foreach \x/\lab in {3.5/{t = 0, stop asserted},3.9/{one bus tick, 0.1 ms},6.8/{cross-arm budget: 3 ms},12.6/{system-wide budget: 500 ms}}{
  \draw[protoblue,line width=0.8pt,dashed] (\x,0.6) -- (\x,-5.2);}
\node[mmlabel,anchor=south,font=\tiny\sffamily\bfseries] at (3.5,0.65)  {t = 0};
\node[mmlabel,anchor=south,font=\tiny\sffamily\bfseries] at (4.9,1.25)  {one bus tick, 0.1 ms};
\node[mmlabel,anchor=south,font=\tiny\sffamily\bfseries] at (6.8,0.65)  {cross-arm budget, 3 ms};
\node[mmlabel,anchor=south,font=\tiny\sffamily\bfseries] at (12.6,0.65) {system-wide budget, 500 ms};
\fill[pablue2,opacity=0.45] (3.5,-6.1) rectangle (6.8,-5.6);
\node[font=\tiny\sffamily,anchor=west] at (3.6,-5.85) {cross-arm halt};
\fill[pagrayl,opacity=0.75] (6.8,-6.1) rectangle (12.6,-5.6);
\node[font=\tiny\sffamily,anchor=west] at (6.9,-5.85) {remaining system-wide margin};
\node[umlnote,text width=64mm,anchor=north west] at (-3.6,-7.6)
  {Three budget lines, all protocol-specified, none aspirational: heartbeat period 0.1 ms, cross-arm emergency stop at most 3 ms, system-wide at most 500 ms. The gap between the cross-arm halt and the system-wide budget is deliberate margin.};
\node[umlnote,text width=64mm,anchor=north west] at (5.2,-7.6)
  {For scale, a human blink is roughly 100 to 150 ms. Every arm is held before a blink completes, and the system-wide guarantee is under four blinks. Three parties may press stop, and any one of them is sufficient.};
\end{pafig}
\vspace{-0.7cm}
\figcaption{Figure 18. The emergency-stop chain on a logarithmic millisecond axis, from the instant the\\
surgeon presses stop through the bus broadcast and arm deceleration to zero tip force, with\\
all three protocol budget lines marked across every one of the timelines.}
\end{pafloat}

\begin{pafloat}
\begin{pafig}[mermaid-type: sequenceDiagram]
\node[mmactor,text width=24mm] (S) at (0,0)    {Sensor stack\\640 ch, 10 kHz};
\node[mmactor,text width=24mm] (L) at (4.1,0)  {On-premises LLM\\advisory only};
\node[mmactor,text width=24mm] (G) at (8.2,0)  {Safety gate\\deterministic};
\node[mmactor,text width=24mm] (H) at (12.3,0) {Operating surgeon\\sole authority};
\node[mmactor,text width=24mm] (R) at (16.4,0) {Robot arms\\8, force-capped};
\foreach \n in {S,L,G,H,R}{\draw[mmlife] (\n.south) -- ++(0,-15.4);}
\fill[pablue1,draw=protoblack,line width=0.4pt] (3.97,-1.0) rectangle (4.23,-3.4);
\fill[pablue1,draw=protoblack,line width=0.4pt] (8.07,-3.4) rectangle (8.33,-9.0);
\fill[pablue1,draw=protoblack,line width=0.4pt] (12.17,-9.0) rectangle (12.43,-11.4);
\fill[pablue1,draw=protoblack,line width=0.4pt] (16.27,-11.4) rectangle (16.53,-13.4);
\draw[mmmsg] (S|-0,-1.0) -- node[above]{1: current field state, force vector, vessel proximity} (L|-0,-1.0);
\draw[mmmsg] (4.23,-1.8) to[out=25,in=-25,looseness=6] node[right,pos=0.5,xshift=2pt]{2: propose next step, no actuation path exists} (4.23,-2.6);
\draw[mmret] (L|-0,-3.4) -- node[above]{3: candidate plan, confidence, rationale} (G|-0,-3.4);
\draw[mmmsg] (8.33,-4.2) to[out=25,in=-25,looseness=6] node[right,pos=0.5,xshift=2pt]{4: check tip force $\leq 3$ N per arm, $\leq 18$ N summed} (8.33,-5.0);
\draw[mmmsg] (8.33,-5.4) to[out=25,in=-25,looseness=6] node[right,pos=0.5,xshift=2pt]{5: check no-fly envelope, SMV and portal vein} (8.33,-6.2);
\draw[protogray,line width=0.5pt] (2.4,-6.5) rectangle (18.2,-9.4);
\node[font=\tiny\sffamily\bfseries,text=protogray,anchor=north west] at (2.5,-6.5) {alt};
\draw[protogray,line width=0.4pt,dashed] (2.4,-8.0) -- (18.2,-8.0);
\node[font=\tiny\sffamily\itshape,anchor=north west] at (3.3,-6.52) {[any limit violated]};
\node[font=\tiny\sffamily\itshape,anchor=north west] at (3.3,-8.02) {[all limits satisfied]};
\draw[mmret] (G|-0,-7.5) -- node[above]{6: rejected, reason logged, never shown as executable} (L|-0,-7.5);
\draw[mmmsg] (G|-0,-9.0) -- node[above]{7: presented for approval, with confidence and rationale} (H|-0,-9.0);
\draw[mmmsg] (12.43,-9.9) to[out=25,in=-25,looseness=6] node[right,pos=0.5,xshift=2pt]{8: accept, modify, or decline; declining needs no justification} (12.43,-10.7);
\draw[mmmsg] (H|-0,-11.4) -- node[above]{9: approved motion command, signed by the surgeon} (R|-0,-11.4);
\draw[mmmsg] (16.53,-12.0) to[out=25,in=-25,looseness=6] node[right,pos=0.5,xshift=2pt]{10: execute inside the envelope, heartbeat maintained} (16.53,-12.8);
\draw[mmret] (R|-0,-13.4) -- node[above]{11: resulting state} (S|-0,-13.4);
\node[umlnote,text width=70mm,anchor=north west] at (0,-16.0)
  {The surgeon may press stop at any instant. A cross-arm halt completes in at most 3 ms and a system-wide halt in at most 500 ms, which is Figure 18.};
\node[umlnote,text width=70mm,anchor=north west] at (9.4,-16.0)
  {There is no arrow from the model to the arms. Not a throttled one, not a supervised one, not a fallback one. That absence is the safety argument.};
\end{pafig}
\vspace{-0.7cm}
\figcaption{Figure 19. One operative step message by message across five lifelines, with the deterministic\\
gate sitting between advice and approval, and no arrow at all from the model to the arms,\\
which is the safety argument this figure makes by an absence rather than words.}
\end{pafloat}

\subsection{The vascular no-fly corridor and the force envelope}

A force cap does not prevent an arm entering a vascular corridor, and a corridor
does not cap the force applied inside permitted tissue. Neither limit is
sufficient alone, and the pairing is the safety property. Figure 20 draws them
as sibling containers inside one envelope for that reason, with the
deterministic gate as the single element that reads both.

The corridor is 4 mm either side of the superior mesenteric and portal veins,
and the approach angle to the superior mesenteric artery is constrained
separately. Proximity is tested at the heartbeat rate, so a drifting instrument
is caught within one tick rather than at the end of a motion.

\begin{pafloat}
\begin{pafig}[d2-type: container with an embedded measurement grid]
\node[d2title,anchor=south west] at (-7.6,3.15) {operative.envelope};
\node[d2title2,anchor=south west] at (-7.2,2.55) {envelope.force\_limits};
\node[d2title2,anchor=south west] at (0.9,2.55) {envelope.no\_fly\_geometry};
\begin{scope}[shift={(-7.0,0)}]
  \foreach \i/\lab/\cap/\obs/\capl/\obsl in {%
    0/{Tip force, one arm}/{3.00}/{2.10}/{3.0 N}/{2.1 N},
    1/{Tip force, summed}/{3.00}/{2.07}/{18.0 N}/{12.4 N},
    2/{Positional error}/{3.00}/{1.95}/{2.0 mm}/{1.3 mm},
    3/{Force reproducibility}/{3.00}/{1.86}/{0.5 N}/{0.31 N},
    4/{Heartbeat period}/{3.00}/{3.00}/{0.1 ms}/{0.1 ms}}{
    \node[font=\tiny\sffamily,anchor=east] at (-0.15,{1.7-\i*0.82}) {\lab};
    \fill[pablue1] (0,{1.55-\i*0.82}) rectangle (\obs,{1.85-\i*0.82});
    \draw[protoblack,line width=0.3pt] (0,{1.55-\i*0.82}) rectangle (\obs,{1.85-\i*0.82});
    \draw[protoblue,line width=0.9pt,dashed] (\cap,{1.48-\i*0.82}) -- (\cap,{1.92-\i*0.82});
    \node[font=\tiny\sffamily,anchor=west] at (3.15,{1.7-\i*0.82}) {cap \capl, worst \obsl};}
  \node[font=\tiny\sffamily\itshape,text=protogray,anchor=north west] at (-2.9,-2.6)
    {Bar: worst value observed across the Phase 0 simulation set.\\Dashed rule: the protocol cap. The gap is the headroom.};
\end{scope}
\begin{scope}[shift={(4.6,0)}]
  \fill[pagraym,opacity=0.55] (-0.45,-2.1) rectangle (0.45,2.1);
  \fill[pagraym,opacity=0.55] (1.55,-2.1) rectangle (2.45,2.1);
  \fill[pablue1] (-0.07,-2.1) rectangle (0.07,2.1);
  \fill[pablue1] (1.93,-2.1) rectangle (2.07,2.1);
  \node[font=\tiny\sffamily,anchor=south,text=protoblue] at (0,2.15) {SMV};
  \node[font=\tiny\sffamily,anchor=south,text=protoblue] at (2.0,2.15) {portal vein};
  \node[font=\tiny\sffamily,anchor=north,text=protogray,align=center] at (1.0,-2.2) {hatched band: 4 mm no-fly\\corridor either side};
  \foreach \x/\y in {-1.2/1.5,-1.0/0.4,-1.3/-0.8,-0.9/-1.7,3.2/1.4,3.0/0.3,3.4/-0.9,3.1/-1.8}{
    \fill[protoblue] (\x,\y) circle (0.075);}
  \node[font=\tiny\sffamily,anchor=east,align=right] at (-1.5,0.0) {eight arm tips,\\none inside a corridor};
  \node[font=\tiny\sffamily,anchor=west,align=left] at (3.6,0.0) {approach angle to the SMA\\constrained separately};
\end{scope}
\node[d2dark,text width=140mm,minimum height=9mm] at (0.3,-4.2)
  {Invariant: a motion plan that would violate any row on the left, or enter any corridor on the right, is rejected by a deterministic gate before it is shown to the surgeon as executable.};
\begin{scope}[on background layer]
\node[d2cont2,fit={(-7.4,2.4) (-2.6,-3.4)}] {};
\node[d2cont2,fit={(2.0,2.4) (9.4,-3.4)}] {};
\node[d2cont,fit={(-7.8,2.9) (9.8,-4.9)}] {};
\end{scope}
\end{pafig}
\vspace{-0.7cm}
\figcaption{Figure 20. The operative envelope drawn as two sibling containers, the force limits with\\
their observed headroom on the left and the vascular no-fly geometry on the right, joined\\
by the deterministic gate that rejects any plan violating either of them.}
\end{pafloat}

\subsection{What could fail, and what catches it}

Figure 21 is a fault tree, and the gate at its apex is an AND gate. Harm reaches
the participant only if all four intermediate events occur together: a
mechanical or control fault, a failure of the deterministic gate to reject it,
no human intervening in time, and the fault not being visible in the telemetry.
Each of those four is itself decomposed, and each has a named barrier.

Four barriers are intact. The force caps are hardware-limited. The no-fly
corridor is enforced by the gate. Surgeon approval is required for every motion,
with a second station and an independent safety monitor able to stop. The 10 kHz
heartbeat halts all arms within 3 ms without a human if it is lost.

One is not, and the figure says so in medium-dark gray: registration drift
beyond 2 mm is caught by the deterministic gate and by nothing else. It is a
single point of failure in the barrier set. Naming it is what a fault tree is
for, and suppressing it would have made the other four less credible.

\begin{pafloat}
\begin{pafig}[graphviz-type: digraph, fault tree with AND and OR gates]
\node[gvboxd,text width=40mm,minimum height=9mm] (top) at (0,0) {\textbf{TOP EVENT}\\harm reaches the participant};
\umlgateand{0}{-1.5}
\node[gvboxs,text width=26mm] (e1) at (-6.6,-3.4) {A mechanical or control fault occurs};
\node[gvboxs,text width=26mm] (e2) at (-2.2,-3.4) {The deterministic gate fails to reject it};
\node[gvboxs,text width=26mm] (e3) at (2.2,-3.4)  {No human intervenes in time};
\node[gvboxg,text width=26mm] (e4) at (6.6,-3.4)  {The fault is not visible in the telemetry};
\umlgateor{-6.6}{-4.8}
\umlgateand{-2.2}{-4.8}
\umlgateand{2.2}{-4.8}
\umlgateor{6.6}{-4.8}
\node[gvcircle] (b1) at (-8.4,-7.0) {arm servo over-travel};
\node[gvcircle] (b2) at (-6.6,-7.0) {instrument fracture};
\node[gvcircle] (b3) at (-4.8,-7.0) {vision loss or occlusion};
\node[gvcircle] (b4) at (-3.1,-7.0) {model proposes an unsafe plan};
\node[gvcircle] (b5) at (-1.3,-7.0) {force sensor out of calibration};
\node[gvcircle,fill=pagrayd] (b6) at (0.5,-7.0) {registration drift $> 2$ mm};
\node[gvcircle] (b7) at (2.4,-7.0) {surgeon does not see the alert};
\node[gvcircle] (b8) at (4.2,-7.0) {safety monitor absent};
\node[gvcircle,fill=pagraym] (b9) at (6.0,-7.0) {stop path unavailable};
\node[gvcircle] (b10) at (7.8,-7.0) {heartbeat loss undetected};
\draw[gvedge] (top) -- (0,-1.25);
\foreach \n in {e1,e2,e3,e4}{\draw[gvedge] (0,-1.78) -- (\n.north);}
\foreach \p/\c in {-6.6/b1,-6.6/b2,-6.6/b3,-2.2/b4,-2.2/b5,-2.2/b6,2.2/b7,2.2/b8,2.2/b9,6.6/b10}{
  \draw[gvedge] (\p,-5.08) -- (\c.north);}
\foreach \p in {-6.6,-2.2,2.2,6.6}{\draw[gvedge] (\p,-3.72) -- (\p,-4.54);}
\node[gvboxs,text width=44mm,anchor=north west] (r1) at (-10.4,-8.6) {\textbf{BARRIER} tip force $\leq 3$ N per arm and $\leq 18$ N summed, hardware limited};
\node[gvboxs,text width=44mm,anchor=north west] (r2) at (-4.6,-8.6) {\textbf{BARRIER} no-fly corridor 4 mm around the SMV and portal vein, gate rejects};
\node[gvboxs,text width=44mm,anchor=north west] (r3) at (1.2,-8.6)  {\textbf{BARRIER} surgeon approval required for every motion, plus a second station and the ISM};
\node[gvboxs,text width=44mm,anchor=north west] (r4) at (7.0,-8.6)  {\textbf{BARRIER} 10 kHz heartbeat; loss halts all arms in $\leq 3$ ms without a human};
\node[gvbox,fill=pagrayd,text width=52mm,anchor=north west] (r5) at (-2.4,-10.9) {\textbf{SINGLE BARRIER ONLY} registration drift is caught by the deterministic gate alone; stated here as a gap rather than omitted};
\draw[gvedged] (r1.north) -- (-6.6,-5.3);
\draw[gvedged] (r2.north) -- (-2.2,-5.3);
\draw[gvedged] (r3.north) -- (2.2,-5.3);
\draw[gvedged] (r4.north) -- (6.6,-5.3);
\draw[gvedge] (r5.north west) to[out=120,in=-90,looseness=0.9] (b6.south);
\node[font=\scriptsize\sffamily\itshape,text=protogray,anchor=north west,align=left]
  at (-10.4,-12.2) {The apex is an AND gate: all four intermediate events must occur together for harm to reach the participant. Four barriers are\\intact and one is not, and the one that is not is drawn in medium-dark gray rather than left out of the figure.};
\end{pafig}
\vspace{-0.7cm}
\figcaption{Figure 21. The hazard fault tree with the top event behind an AND gate, so that harm requires\\
all four intermediate events together, and with the single-barrier gap at registration\\
drift marked in medium-dark gray rather than being quietly left out of it.}
\end{pafloat}

\subsection{Conversion to open surgery}

Conversion from the robotic approach to an open operation is a recorded outcome
of the procedure, not a failure of the study, of the device, or of the
participant. It does not end participation, it does not remove the participant
from any analysis population, and the decision belongs to the operating surgeon
alone and is not subject to review during the procedure by anyone else.

An open tray and backup instruments are in the room for every procedure, and the
conversion sequence is rehearsed as part of the pre-procedure checklist. In the
contemporaneous nationwide series of 1000 robotic pancreatoduodenectomies,
conversion was most frequent in the earliest learning-curve phase and fell with
institutional experience \cite{DutchCohort2025}; this study enrols only at sites
already past the third phase of that curve, which is stated in \S 10.2 with the
numbers attached.

If a conversion occurs, the participant is told before they leave hospital, the
reason is recorded in the audit chain, and the event is reported to the DSMB at
the next cohort boundary. None of those three is discretionary.

\subsection{The drug, paused and restarted around the operation}

Daraxonrasib is given before the operation, paused across it, and restarted
afterwards on an advised schedule. Each of those three phases exists for a
reason a participant can check.

\textbf{Before.} The preoperative course runs through the screening and baseline
window at the assigned dose level. Its purpose in this study is not to shrink
the tumour to resectability; that is not what a Phase 1 safety study can claim.
Its purpose is to establish tolerance at the assigned level in the individual
participant before the operation adds its own physiological load.

\textbf{The pause.} The drug is held before surgery for a protocol-specified
interval. The reason is anastomotic healing: three new connections are made
during a Whipple, a pancreaticojejunostomy, a hepaticojejunostomy, and a
gastrojejunostomy or duodenojejunostomy, and the first of those three is the one
whose failure defines the dominant early complication of the operation
\cite{Bassi2017}. Holding a systemic agent across the healing window is the
conservative choice, and this protocol takes it.

\textbf{The restart.} The drug is restarted at T+7, T+14, or T+21 days. Which of
the three is advised depends on two measured quantities: the ISGPS fistula grade
at the time of assessment, and the drug trough. The advisory is produced by the
same on-premises model that advises intra-operatively, it is reviewed by the Site
Principal Investigator, and it is presented to the participant as an advisory
rather than an instruction. The participant may decline it, and declining does
not end participation or change any other part of the schedule.

\vspace{0.30em}

\begin{xltabular}{\textwidth}{L{2.2cm} L{4.2cm} L{3.0cm} Y}
\toprule
\textbf{Restart window} & \textbf{Advised when} & \textbf{Who decides} & \textbf{What you are told} \\
\midrule
\endfirsthead
\multicolumn{4}{@{}l}{\footnotesize\itshape Table continued from the previous page.}\\
\toprule
\textbf{Restart window} & \textbf{Advised when} & \textbf{Who decides} & \textbf{What you are told} \\
\midrule
\endhead
\midrule
\multicolumn{4}{r@{}}{\footnotesize\itshape Table continued on the next page.}\\
\endfoot
\bottomrule
\endlastfoot
T+7 days & No fistula, or ISGPS biochemical leak only, and trough in range & Advisory from the model, review by the Site PI, decision by you & The grade, the trough, and why the earliest window was advised \\
T+14 days & ISGPS grade B fistula resolving, or trough below range at T+7 & Same & What changed since T+7, and what would move it to T+21 \\
T+21 days & ISGPS grade B fistula not yet resolving, or any grade C & Same & Why the drug is being held longer and what the plan is if it does not resolve \\
Not restarted & Grade C fistula unresolved at T+21, or a dose-limiting toxicity & Site PI, with DSMB notification & That surgical follow-up and every safety assessment continue unchanged \\
\end{xltabular}

\vspace{0.30em}

\subsection{The three connections, and why the first one matters most}

A pancreaticoduodenectomy removes the head of the pancreas, the duodenum, the
gallbladder, the distal bile duct, and, in the classical form, the gastric
antrum. Reconstruction then rebuilds three connections in sequence.

The pancreaticojejunostomy joins the remnant pancreas to the small bowel. It is
the connection whose leak produces a postoperative pancreatic fistula, and the
ISGPS grading of that leak, biochemical, grade B, or grade C, is the secondary
endpoint that carries the most weight in the early window \cite{Bassi2017}. It is
also the connection the force caps most directly protect: soft, friable
pancreatic tissue is the tissue least tolerant of traction, and the 3 N per-arm
cap exists for it.

The hepaticojejunostomy joins the bile duct to the bowel. Its leak rate is far
lower and its consequences are usually manageable with drainage.

The gastrojejunostomy or duodenojejunostomy restores gastric outflow. Delayed
gastric emptying after it is common, uncomfortable, and rarely dangerous, and
participants should expect to be told about it because it is the complication
most likely to lengthen their inpatient stay without being serious.

The order matters to the participant for one practical reason: the drain that
stays in for days after the operation is there to detect a leak from the first
of the three connections, and the drain fluid check described in \S 9.1 is what
converts that drain into the ISGPS grade.

\subsection{The day of the operation, hour by hour}

The protocol describes the operation as a procedure. The participant
experiences it as a day, and the two descriptions do not overlap much. What
follows is the day, with the study-specific steps marked, so that a participant
can tell which parts are the operation they would have had anyway and which
parts exist because this is a study.

\vspace{0.30em}

\begin{xltabular}{\textwidth}{L{2.5cm} C{1.5cm} Y}
\toprule
\textbf{Roughly when} & \textbf{Study-specific} & \textbf{What is happening, and where you are} \\
\midrule
\endfirsthead
\multicolumn{3}{@{}l}{\footnotesize\itshape Table continued from the previous page.}\\
\toprule
\textbf{Roughly when} & \textbf{Study-specific} & \textbf{What is happening, and where you are} \\
\midrule
\endhead
\midrule
\multicolumn{3}{r@{}}{\footnotesize\itshape Table continued on the next page.}\\
\endfoot
\bottomrule
\endlastfoot
Day -1, afternoon & Partly & Baseline visit. Bloods, performance status, and the Phase 0 gate sign-off is shown to you. Daraxonrasib is held from this point \\
Day 0, before dawn & No & Admission, marking, and the anaesthetic conversation. Nothing about the platform changes this \\
About 60 min before & Yes & Room readiness check: heartbeat bus verified at 10 kHz, emergency-stop latency measured and logged for your case, no-fly geometry loaded from your own imaging \\
Induction & No & General anaesthesia, lines, and positioning. You are asleep before the arms are brought to the table \\
First 30 to 60 min & No & Access, staging laparoscopy, and the decision that the tumour is still resectable. If it is not, the operation stops here and no platform step occurs \\
The resection & Yes & Docking of the eight arms, then the resection itself, step by step. Every motion is proposed, gated, and signed before it executes \\
The reconstruction & Partly & The three reconnections. The pancreatic anastomosis is the one \S 7.9 discusses; the surgeon may complete any of the three by hand \\
Closure & No & Drains placed, closure, and the count. Standard care \\
Recovery, hours 0 to 4 & No & Post-anaesthesia care unit. Your family is told when you arrive there \\
Ward, day 0 evening & Yes & The audit trail for your case is sealed and hashed. Your emergency-stop and force-envelope numbers exist from this point and may be requested \\
\end{xltabular}

\vspace{0.30em}

\noindent Two things in that table are worth pulling out. The room readiness
check is the only study-specific step before induction, and it happens whether
or not it finds a problem; if it finds one, the operation proceeds without the
platform. And the audit trail is sealed the same evening, which is why the
extract in \S 11.5 can be requested from the following day rather than at study
closure.

\subsection{Who is in the room, and what each of them does}

A participant asked to consent to an eight-arm platform reasonably wants to know
whether the room is more crowded than it would otherwise be, and whether the
additional people are clinical or technical.

The clinical team is unchanged from a conventional robotic Whipple: the
operating surgeon at the console, an assistant surgeon at the table, the
anaesthetist, the scrub practitioner, and the circulating practitioner. Those
five make every clinical decision, and the operating surgeon is the one named in
the participant's consent form.

Two roles are added by the study. An independent safety monitor watches the gate
decisions and the force telemetry and holds an unconditional authority to halt,
which they may exercise without agreeing it with the surgeon first. A platform
engineer maintains the stack and has no clinical role, no console access, and no
authority over any motion. The distinction is deliberate and is visible in the
audit trail: no motion in the record can carry an engineer's identity, and a
record that did would be a protocol violation rather than an unusual event.

The safety monitor is the role most often misread, so it is worth being exact.
They are not a second surgeon and they do not offer an opinion on the operative
plan. Their entire function is to stop the platform when a limit is approached
or a gate behaves unexpectedly, and their stop does not require the surgeon's
agreement. A participant who wants to know whether anyone in the room can
override enthusiasm has a specific answer: yes, and that person's only job is
to be willing to.

\subsection{What is different from a conventional robotic Whipple}

Most of what happens is not new. The comparison below separates what the
platform changes from what it does not, because a participant weighing the study
is weighing the difference and not the whole operation.

\vspace{0.30em}

\begin{xltabular}{\textwidth}{L{4.3cm} Y}
\toprule
\textbf{Aspect} & \textbf{Conventional robotic Whipple, and what this study changes} \\
\midrule
\endfirsthead
\multicolumn{2}{@{}l}{\footnotesize\itshape Table continued from the previous page.}\\
\toprule
\textbf{Aspect} & \textbf{Conventional robotic Whipple, and what this study changes} \\
\midrule
\endhead
\midrule
\multicolumn{2}{r@{}}{\footnotesize\itshape Table continued on the next page.}\\
\endfoot
\bottomrule
\endlastfoot
The operation performed & Identical. Same resection, same specimen, same three reconnections \\
Who moves the instruments & Identical. The surgeon moves them from a console; no motion is autonomous in either case \\
Number of arms & Eight rather than four. The additional arms hold retraction and camera positions a human assistant would otherwise hold \\
Force at the tip & New. A hardware cap at 3 N per arm and 18 N summed, which a conventional platform does not enforce \\
Vascular protection & New. A 4 mm no-fly corridor computed from your own imaging, enforced by the gate rather than by vigilance alone \\
Step planning & New. The model proposes the next step with a rationale; the surgeon accepts, edits, or declines it \\
Stopping & Stronger. A measured 3 ms cross-arm and 500 ms system-wide budget, verified for your case before induction \\
Record & New. Every advisory, gate decision, and approval is hash-chained and available to you \\
Operative time & Longer during a learning phase, and \S 7.7 states by how much and why that matters \\
Anaesthesia and recovery & Identical. Nothing in the platform changes the anaesthetic or the ward pathway \\
\end{xltabular}

\vspace{0.30em}

\noindent Read down the table, the platform adds four constraints, one record,
and one advisory layer, and changes nothing about what is cut or how it is
sewn. That is a narrower claim than the phrase Physical AI usually suggests, and
it is the accurate one.
